\documentclass[a4paper]{report}
% Shared preamble for the English probability and statistics documents.
% Based on the original course preamble. Compile with pdfLaTeX.
\usepackage[top=1.5cm,bottom=2cm,left=2cm,right=2cm]{geometry}
\usepackage[T1]{fontenc}
\usepackage{lmodern}
\usepackage[utf8]{inputenc}
\usepackage[english]{babel}

\usepackage{amsmath}
\usepackage{amssymb}
\usepackage{amsthm}
\usepackage{mathrsfs}
\usepackage{mathtools}
\usepackage{etoolbox}
\usepackage{tikz}
\usetikzlibrary{shapes.multipart,arrows.meta}
\usepackage{xcolor}
\usepackage{enumerate}
\usepackage[shortlabels]{enumitem}
\usepackage{hyperref}
\usepackage{version}
\usepackage{bbold}
\usepackage{fontawesome5}
\usepackage{fancyhdr}
\lfoot{}
\usepackage{natbib}
\usepackage{graphicx}
\usepackage{pdfpages}
\usepackage{authblk}
    
\usepackage{blindtext}
\usepackage{multicol}
\usepackage{collcell}
\usepackage{wrapfig}
\usepackage{blkarray}
\usepackage{cancel}
\usepackage{appendix}
\usepackage{float}
\makeatletter
\input{tex/assets/slashbox.sty}
\makeatother
\hypersetup{unicode,colorlinks,citecolor=red,filecolor=black,
  linkcolor=red,urlcolor=cyan}
\setlength{\columnsep}{1cm}
\pagenumbering{gobble}

\def\numtext#1{%
  \ifcase#1\or one\or two\or three\or four\or five\or six\or seven\or eight\or nine\or ten\or
  eleven\or twelve\or thirteen\or fourteen\or fifteen\or sixteen\or seventeen\or eighteen\or nineteen\or twenty\or Lots\fi}

\newcommand{\Byte}[1]{%
  \def\nodecontents{}%
  \foreach[count=\l] \k in {#1} {\xappto{\nodecontents}{\noexpand\nodepart{\numtext{\l}}\k}}%
  \node[rectangle split,rectangle split parts=8,rectangle split horizontal,draw] at (2,2){\nodecontents};}

\newcommand{\D}{{\mathrm{d}}}
\newcommand{\cE}{{\mathcal{E}}}
\newcommand{\cG}{{\mathcal{G}}}
\newcommand{\cX}{{\mathcal{X}}}
\newcommand{\cY}{{\mathcal{Y}}}
\newcommand{\cB}{{\mathcal{B}}}
\newcommand{\cF}{{\mathcal{F}}}
\newcommand{\cO}{{\mathcal{O}}}
\newcommand{\cS}{{\mathcal{S}}}
\newcommand{\cT}{{\mathcal{T}}}
\newcommand{\given}{\, \big| \,}
\newcommand{\cL}{{\mathscr{L}}}
\newcommand{\cD}{{\mathcal{D}}}
\newcommand{\cA}{{\mathscr{A}}}
\newcommand{\cP}{{\mathscr{P}}}
\newcommand{\cR}{{\mathscr{R}}}
\newcommand{\A}{{\mathcal A}}
\newcommand{\one}{\mathbb{1}}
\newcommand{\mathbbm}[1]{\mathbb{#1}}

\providecommand{\C}{}
\renewcommand{\C}{\mathbb C}
\newcommand{\R}{\mathbb R}
\newcommand{\Q}{\mathbb Q}
\newcommand{\N}{\mathbb N}
\newcommand{\T}{\mathbb T}
\newcommand{\F}{\mathbb F}
\renewcommand{\P}{\mathbb P}
\newcommand{\Z}{\mathbb Z}
\newcommand{\E}{\mathbb {E}}
\newcommand{\Pbb}[1]{\P\left( #1\right)}
\newcommand{\Ebb}[1]{\E\left[ #1\right]}

\newcommand{\lb}{\left (}
\newcommand{\rb}{\right )}
\newcommand{\lbb}{\left [}
\newcommand{\rbb}{\right ]}
\newcommand{\labs}{\left |}
\newcommand{\rabs}{\right |}
\newcommand{\lbrb}[1]{\lb #1 \rb}
\newcommand{\lbbrbb}[1]{\lbb#1\rbb}
\newcommand{\labsrabs}[1]{\labs#1\rabs}
\newcommand{\lbbrb}[1]{\lbb#1\rb}
\newcommand{\lbrbb}[1]{\lb#1\rbb}
\newcommand{\langlerangle}[1]{\langle#1\rangle}
\newcommand{\lbcurly}{\{}
\newcommand{\rbcurly}{\}}
\newcommand{\curly}[1]{\lbcurly#1\rbcurly}
\newcommand{\abs}{\labsrabs}
\newcommand{\ind}[1]{1_{#1}}
\newcommand{\indep}{\perp \!\!\! \perp}
\renewcommand{\d}{\, \mathrm{d}}

\newtheorem{theorem}{Theorem}
\newtheorem{problemCnt}{Lemma}
\theoremstyle{definition}
\newtheorem{problem}[problemCnt]{Problem}
\newtheorem{problembook}[problemCnt]{\faBook \,  Problem}
\newtheorem{problempencil}[problemCnt]{\faPencil* \,  Problem}
\newtheorem{problemstar}[problemCnt]{\faStar\,  Problem}
\newtheorem{problemstarpencil}[problemCnt]{\faStar \, \faPencil* \,   Problem}
\newtheorem{probleminteresting}[problemCnt]{\faSmile[regular] \,   Problem}
\newtheorem{solution}{Solution}

\newcommand{\fwcell}[1]{\makebox[\arraystretch\normalbaselineskip]{$#1$}}
\newenvironment{enumeratei}{\begin{enumerate}[(i)]}{\end{enumerate}}
\newenvironment{enumeratea}{\begin{enumerate}[(a)]}{\end{enumerate}}
\newenvironment{enumeraten}{\begin{enumerate}[\upshape 1.]}{\end{enumerate}}

% Compatibility helpers needed by a few files in the combined archive.
\AtBeginDocument{\noindent{\footnotesize English translation prepared with ChatGPT (OpenAI, GPT-6), 26 October 2026, for translation of the Bulgarian text into English and LaTeX formatting.
The Bulgarian originals remain the reference; translation may contain errors.}\par\smallskip}
\newcommand{\MissingGraphic}[1]{%
  \begin{center}\fbox{\parbox{0.82\linewidth}{\centering Missing source image: \texttt{#1}}}\end{center}}



\newif\ifshowsolutions
\showsolutionstrue
\begin{document}
\noindent Probability and Statistics, CS Year 3, 2020/2021. Name:\hfill Student no.:\qquad\qquad
\hrulefill
\begin{it}Your grade is 1 + the number of points. Good luck! Work on $(\Omega,\mathcal F,\P)$. Define $\mathbb N_0=\mathbb N\cup\{0\}$. Unless otherwise stated, $n,k,i\in\mathbb N$.\end{it}
\begin{problem}
An attacker sends a phishing email to a lecturer. At a random time $X\sim U(0,1)$ during the next day, the attacker hides all traces but also loses the opportunity to benefit. Thus $f_X(x)=1$ on $[0,1]$ and 0 elsewhere. The lecturer responds after an independent exponential time with mean two days: $Y\sim\mathrm{Exp}(1/2)$, with $f_Y(y)=\frac12e^{-y/2}$ for $y\geq0$ and 0 elsewhere.
\begin{enumerate}
\item (0.75 points) What is the probability that the attack succeeds? Approximately what fraction of many attempts will succeed?
\ifshowsolutions
\par\textbf{Solution.} The desired probability is
\[p:=\P(X<Y)=\int_0^1\int_0^x\frac12e^{-y/2}\,dy\,dx,\]
which is easily calculated. By the law of large numbers, the successful fraction is $p$; specify the variables and verify the hypotheses.
\fi
\item (0.25 points) What is the probability the attacker benefits given that the lecturer responds after exactly 12 hours?
\ifshowsolutions
\par\textbf{Solution.} We seek $\P(Y<X\mid X=1/2)=\P(Y<1/2)$. The intention was to use a conditional density, but it is actually unnecessary.
\fi
\end{enumerate}
\end{problem}
\begin{problem}
Two faces of a fair die are white, two green and two red. Roll it twice. Let $X$ count white results and $Y$ red results. Find:
\begin{enumerate}
\item (0.5 points) the joint distribution and correlation of $X,Y$. Are they independent?
\item (0.25 points) $\P(X=1\mid Y=1)$;
\item (0.25 points) $\P(X>Y)$.
\end{enumerate}
\ifshowsolutions\par\textbf{Solution in the source:} ---\fi
\end{problem}
\begin{problem}
Daily accident counts in Bulgaria are modelled by independent Poisson variables $X_1,X_2,\ldots$ with $\lambda=10$, so $\P(X_1=k)=e^{-\lambda}\lambda^k/k!$ for $k\in\mathbb N_0$. Set $S_n=X_1+\cdots+X_n$.
\begin{enumerate}
\item (0.25 points) Find $\E S_n$ and $DS_n$.
\ifshowsolutions\par\textbf{Solution.} $\E S_n=DS_n=n\lambda$, since the $X_i$ are independent and $\E X_i=DX_i=\lambda$.\fi
\item (0.75 points) If $T$ is the total number of serious accidents in 2021, show $\P(3529\leq T\leq3771)\approx0.954$ and carefully justify the approximation.
\ifshowsolutions\par\textbf{Solution.} Use the CLT, explicitly checking its hypotheses, in particular $\E[X_i^2]<\infty$.\fi
\end{enumerate}
\end{problem}
\begin{problem}
Each of $n>2$ people tosses a fair coin. A person wins if their result differs from everyone else's. If there is no winner, all toss again.
\begin{enumerate}
\item (0.25 points) Let $X$ count rounds until a winner. Find its expectation and variance.
\ifshowsolutions
\par\textbf{Solution.} A winner requires one head and all other tails, or the reverse. This has probability $p_n=n/2^{n-1}$ (why?). Hence $X\sim\mathrm{Ge}(p_n)$, $\E X=1/p_n$ and $DX=1/p_n^2$.
\fi
\item (0.75 points) After the first winner leaves, continue with the remaining $n-1$ players, then $n-2$, and so on until two remain. What is the expected number of rounds? If winner $k$ receives $100(n-k)$, what is the most you would pay to enter?
\ifshowsolutions
\par\textbf{Solution.} Explain why the expected number is $1/p_n+1/p_{n-1}+\cdots+1/p_3$. By symmetry all players have the same expected winnings, and the total paid is $100[(n-1)+(n-2)+\cdots+2]$, so divide this total by $n$.
\par Although not required, the total simplifies to $100[n(n-1)/2-1]$. The \href{https://math.stackexchange.com/questions/888354/asymptotic-behavior-of-sum-limits-k-1n-frac2kk}{asymptotic behaviour of the other sum} is nontrivial and equals $2^n/n$.
\fi
\end{enumerate}
\end{problem}
\begin{problem}
A fan collects an album of $n=500$ different Euro 2020 stickers (in 2021), buying random stickers. Let $X$ be the number needed to collect all types.
\begin{enumerate}
\item (0.5 points) Let $X_1=1$, and for $i\geq2$ let $X_i$ count purchases after obtaining the $(i-1)$st distinct sticker through obtaining the $i$th, inclusive. For example, the sequence $\mathbf5,\mathbf{12},5,\mathbf{111},12,111,5,\mathbf{13}$ gives $X_1=1,X_2=1,X_3=2,X_4=4$. What is the distribution of $X_i$? Find $\E X$.
\ifshowsolutions
\par\textbf{Solution.} $X_2\sim\mathrm{Ge}(499/500)$, since a new type appears with probability $499/500$; otherwise it is the initial type. Then $X_3\sim\mathrm{Ge}(498/500)$, and so on. Complete the argument.
\fi
\item (0.5 points) Estimate the probability that at least $2n$ purchases are needed.
\ifshowsolutions
\par\textbf{Solution.} Chebyshev's inequality. This is the well-known \href{https://en.wikipedia.org/wiki/Coupon_collector\%27s_problem}{coupon collector's problem}, also appearing in bonus homework 2. Problems 4 and 5 are in fact very similar.
\fi
\end{enumerate}
\end{problem}
\end{document}
